\documentclass[a4paper,12pt]{article}
\usepackage{subcaption}
\usepackage[labelformat=simple,labelsep=period,skip=3pt]{caption}
\usepackage{graphicx}
%%Mit diesem Paket werden Grafiken ausgeblendet, damit die Kompelierung schneller erfolgt
%\usepackage[draft]{graphicx}
%%Für die Einbindung von Grafiken
\usepackage{epsfig, graphics}
%% Footnotes oder tables
\usepackage[flushleft]{threeparttable}
%%Für das Seitenformat.
\usepackage{fullpage}
%%Definition der Seitengrösse
\setlength{\textwidth}{16 true cm}
\setlength{\textheight}{24.5 true cm}
\oddsidemargin  0 cm
\evensidemargin 0 cm
\topmargin      0 cm
%%Für den Zeilenabstand: \doublespacing, etc.
\usepackage{setspace}
%%Für die Schriftart
\usepackage{lmodern}
\usepackage[T1]{fontenc}
%%Für Tabellen (midrule, etc)
\usepackage{booktabs,longtable}
%\usepackage[format=hang,justification=centering,singlelinecheck=off]{caption}
%%Für das eurosymbol
\usepackage{eurosym}
%%Für das Literaturverzeichnis
\usepackage{natbib}
\usepackage{cite}
%%Für Textfarbefc
\usepackage{color}
%Def.Table-Notes
\def\fussnoten{\vskip2pt\footnotesize}
%Gliederung in PDf
\usepackage{hyperref}
%Figurenotes
\usepackage[capposition=top]{floatrow}
%Für Textmarkierungen
\usepackage{xcolor,soul}
% for diagrams:
\usepackage{tikz}
\usetikzlibrary{arrows,positioning}
% for floatbarriers
%\usepackage{placeins}
%%Confining those floaty bastards to their section
\usepackage[section]{placeins}
% für align
\usepackage{amsmath}
%\usepackage{bbm}
\newcommand{\mathbbm}[1]{\text{\usefont{U}{bbm}{m}{n}#1}} % from mathbbm.sty
% for tables
\usepackage{tabularx}
\usepackage{array}
\newcolumntype{Y}{>{\centering\arraybackslash}X}
\newcommand{\sym}[1]{\rlap{#1}}

\title{{ECRMALE: Take Home Exam}\vskip1.3cm}

\date{Version: 18 December 2025\\}

\begin{document}

\maketitle

\newpage

\doublespacing

\section{The cannonical model of the returns to schooling}

Table A1 and Figure A1 below are taken from Autor, Goldin and Katz, "Extending the Race Between Education and Technology", AEA Papers \& Proceedings, May 2020. 

\begin{enumerate}

\item[(1.1)] Interpret the theoretical model and its coefficients that Table A1 aims to estimate. 

\item[(1.2)] Interpret the estimates of coefficients reported in Table A1. 

\item[(1.3)] Using Table A1, discuss the time series plotted in Figure A1.

\item[(1.4)] Are there any empirical facts that the canonical model cannot easily explain?

\end{enumerate}

\begin{center}
\begin{figure}
\includegraphics[width=\textwidth]{tableA1.pdf} 
\end{figure}
\end{center}

\begin{center}
\begin{figure}
\includegraphics[width=\textwidth]{figureA1.pdf} 
\end{figure}
\end{center}

\newpage

\section{Automation and labor demand: theory}

\subsection{Environment}

Assume the following competitive economy:

\begin{itemize}
\item Aggregate production of a single final good, $Y$, is given by the following Cobb-Douglas production function:
\[
Y= \exp \left[ \int_{0}^{1} \ln (y(z))dz \right] \tag{2.1}  \label{eq1}
\]
with $y(z)$ the quantity of task $z$ used in the production of $Y$. There is a continuum of tasks spread out over the unit-interval such that $z \in [0,1]$.
\item Each task is produced using a quantity of capital, $k(z)$, or of labor, $l(z)$, according to:
\[
y(z) = 
\begin{cases}
\gamma^{K}(z)k(z) \text{ if } z \in [0,I] \tag{2.2} \\
\gamma^{L}(z)l(z) \text{ if } z \in (I,1] 
\end{cases} \label{eq2}
\]
with $\gamma^{K}(z)$ and $\gamma^{L}(z)$ the task-specific productivity of capital and labor, respectively. Task $z = I \in (0,1)$ is an exogenous task threshold such that all tasks $z \leq I$ can only be produced by capital and all tasks $z > I$ can only be produced by labor.
\item In equilibrium, clearing of capital and labor markets requires that:
\[
\int_{0}^{1} k(z)dz = K \text{ and } \int_{0}^{1} l(z)dz = L  \tag{2.3} \label{eq3}
\] 
with $K$ and $L$, respectively, a given aggregate quantity of capital and labor in the economy.
\item Given the Cobb-Douglas production function in equation (\ref{eq1}), the corresponding expression for the marginal cost of producing $Y$ is given by:
\[
MC = \exp \left[ \int_{0}^{1} \ln(p(z))dz \right] \tag{2.4} \label{eq4}
\]
with $p(z)$ the price or the unit-cost of task $z$.
\end{itemize}

\subsection{Equilibrium}

\begin{enumerate}

\item[(2.1)] Define $R$ as the price of capital and $W$ as the price of labor. That is, $R$ is the price of one unit of $k(z)$ and $W$ is the price of one unit of $l(z)$. For given values of $R$ and $W$, use that marginal revenue products of capital and labor must equal their marginal costs in equilibrium to replace the question marks in the following expression for the unit-cost of task $z$, $p(z)$: \medskip
\[
p(z) = 
\begin{cases}
\textbf{?} \text{ if } z \in [0,I] \tag{2.5} \label{eq5} \\
\textbf{?} \text{ if } z \in (I,1] 
\end{cases}
\]

\item[(2.2)] Given that equation (\ref{eq1}) is a Cobb-Douglas production function using a continuum of tasks, cost shares must be constant and equal across all tasks in equilibrium. In particular, given that tasks are spread out over the unit-interval, we must have that $ \forall z: p(z)y(z)=Y$. Using that $y(z)=Y/p(z)$ together with equations (\ref{eq2}) and (\ref{eq5}), show that:  
\[
k(z)=
\begin{cases}
\frac{Y}{R} \text{ if } z \in [0,I] \tag{2.6} \label{eq6} \\
0 \text{ if } z \in (I,1]
\end{cases}
\]
and
\[
l(z)=
\begin{cases}
0 \text{ if } z \in [0,I]  \tag{2.7} \label{eq7} \\
\frac{Y}{W} \text{ if } z \in (I,1]
\end{cases}
\]
which gives the demand for capital and labor for each task $z$, respectively.

\item[(2.3)] Use equations (\ref{eq3}), (\ref{eq6}) and (\ref{eq7}) to solve for equilibrium expressions for $R$ and $W$.

\item[(2.4)] In equilibrium, the marginal cost of $Y$, given by equation (\ref{eq4}), must equal the price of $Y$ which we denote by $P$. Choosing $Y$ as the numeraire, we must therefore have that $P=MC=1$ in equilibrium. Taking logarithms such that $\ln(P)=\ln(MC)=0$ and using equations (\ref{eq4}) and (\ref{eq5}) together with the equilibrium expressions for $R$ and $W$, re-write $Y$ as the following aggregate Cobb-Douglas production function:
\[
Y = A \left[ \frac{K}{I} \right]^{I} \left[ \frac{L}{1-I} \right]^{1-I} \tag{2.8} \label{eq8}
\]
where $A$ is defined as:
\[
A \equiv \exp \left[ \int_{0}^{I} \ln(\gamma^{K}(z))dz + \int_{I}^{1} \ln(\gamma^{L}(z))dz \right] 
\]

\item[(2.5)] Discuss how the allocation of capital and labor across tasks enters equation (\ref{eq8}).

\end{enumerate}


\subsection{Automation and labor demand}

\begin{enumerate}

\item[(2.6)] Assume that tasks are ranked on the unit-interval such that $\gamma^{K}(z)/ \gamma^{L}(z)$ is decreasing in $z$. That is, capital has a comparative advantage in the production of lower-indexed tasks, and labor has a comparative advantage in the production of higher-indexed tasks. Now assume that technological progress consists of the automation by capital of labor tasks. As examples, think of the automation by industrial robots of car assembly lines or of software automating tasks done by office clerks. In our model, assume that automation is captured by an increase in $I$. Do you find this assumption reasonable?

\item[(2.7)] Taking the logarithm of the expression for $W$ and differentiating, show that the impact of automation on (the logarithm of) $W$ consists of two effects. The first is a negative \textit{direct displacement effect} because the range of tasks done by workers decreases. The second is a \textit{productivity effect} that captures a change in (the logarithm of) average labor productivity.

\item[(2.8)] Taking the logarithm of equation (\ref{eq8}) and differentiating with respect to $I$, proof that the productivity effect will be strictly positive if and only if:
\[
\frac{R}{\gamma^{K}(I)} < \frac{W}{\gamma^{L}(I)}
\]
and interpret this inequality.

\item[(2.9)] Finally, show that automation always decreases the aggregate labor share in the economy, defined by $s \equiv WL/Y$.

\end{enumerate}

\newpage

\section{Automation and labor demand: evidence}

Bessen, Goos, Salomons and van den Berge (2025) estimate the impact of firm-level automation by combining Dutch micro-data with a direct measure of automation expenditures covering all private non-financial sector firms. Using a novel difference-in-differences event-study design leveraging lumpy investment in automation, their Figure 2 copied below finds that automation reduces firm-level employment but not the firm's average wage in companies employing fewer than 500 workers. 

For their analyses, they set up their data in a stacked difference-in-differences design. More specifically, their first cohort of treated firms are firms that have their first automation event in 2003. The event window surrounding this event contains the calendar years from 2000 to 2007. All control firms for the 2003-cohort are those firms that are observed over the same event window (2000 to 2007), and that have their first automation event later than 2007. These firms are ``clean controls'' in the sense that they do not have an automation event during the event window. They repeat this procedure for each cohort of firms. Finally, they stack the cohort-specific datasets so that they line up in terms of event time $\tau \in \{-3,..,4\}$. 

\begin{figure}
\caption*{Figure 2: Firm-level outcomes for automating firms using event timing}\label{firm_event}
\begin{subfigure}[b]{1\textwidth}
	\centering
\caption{Log employment change}
    \includegraphics[width=0.8\textwidth]{firm_stacked_did_weight_emp_500.pdf}
\end{subfigure}
\begin{subfigure}[b]{1\textwidth}
	\centering
\caption{Log daily wage change}
    \includegraphics[width=0.8\textwidth]{firm_stacked_did_weight_wage_500.pdf}
\end{subfigure}
\smallskip 
\begin{minipage}{0.95\textwidth} 
{\footnotesize \textit{Notes}: All automating firms that exist in all years $\tau \in \{-3,..,4\}$. Both models are weighted by firm-level employment size in $\tau=-1$. Whiskers reflect 95\% confidence intervals.\par}
\end{minipage}
\end{figure}

\subsection{The impact of automation on firm-level employment}

\begin{itemize}

\item[(3.1)] To obtain the point estimates plotted in Figure 2, they use their stacked dataset to estimate the following difference-in-differences specification:
\begin{equation} \label{eq_eventstudy}
\ln Y_{jt} = \alpha + \sum_{\tau\neq-1;\tau=-3}^{4}\beta_{\tau}\times I_{\tau} + \sum_{\tau\neq-1;\tau=-3}^{4}\delta_{\tau}\times I_{\tau} \times treat_j +  \eta_j + \theta_t + \varepsilon_{jt} \tag{3.1}
\end{equation}
where $j$ indexes firms by cohort. $I_{t}$ are leads and lags in event time, with $\tau=-1$ as the reference category. Treatment is indicated by the dummy variable $treat_j$ which turns on for firms that have an automation event in the window. $Y_{jt}$ is $j$'s employment or average daily wage and $\eta_j$ are a set of cohort by firm fixed effects and $\theta_t$ a set of calendar year fixed effects. Explain why this estimating equation corresponds to a Difference-in-Differences specification. 

\item[(3.2)] What are the identifying assumptions needed to interpret the coefficients in Figure 2 as causal estimates? Do you think these identifying assumptions are likely to hold? Can you give an example of when they would not hold?

\end{itemize}

\subsection{Difference-in-Differences Event-Study}

\begin{itemize}

\item[(3.3)] If one would estimate equation (\ref{eq_eventstudy}) without first stacking the data as described above, point estimates could be inconsistent (even if the identifying assumptions that you mentioned above would hold). Why?

\item[(3.4)] Why is stacking the data as is done in Bessen et al (2025) a solution to this problem?

\item[(3.5)] Also other solutions have been proposed recently. Briefly describe intuitively one other solution (other than stacking the data as Bessen et al (2025) do).

\end{itemize}

\newpage

\section{Firm wage premia}

\subsection{Estimating AKM models}

Table III, taken from Card, Heinig and Kline (2013), sumarizes estimates of AKM models for different time periodes using German data. Discuss the AKM model as well as the estimates presented in Table III below.  

\begin{figure}
\includegraphics[width=\textwidth]{Table3.png} 
\end{figure}

\subsection{Job differentiation, monopsony, and firm wage premia}

One possible explanation for the existence of firm wage premia is the model in Card, Cardoso, Heinig and Kline (2018).  

\subsubsection{Worker preferences over non-pecuniary amenities}

\begin{itemize}

\item Assume that workers have different preferences over the same workplace amenities. In particular, assume that utility of worker $i$ from working in firm $j$ is given by:
\begin{equation} \label{eq_worker_preferences}
    u_{ij} = \epsilon \ln(w_{j}) + \eta_{ij}
\end{equation}
where $\eta_{ij}$ captures idiosyncratic preferences for working at firm $j$, arising, for example, from non-pecuniary match factors such as distance to work or interactions with co-workers and supervisors. If $\{\eta_{ij}\}$ are independent draws from a type-I Extreme Value distribution, the total number of workers that will be employed at firm $j$ is given by:
\begin{equation} 
    l_{j} = \frac{\exp{(\epsilon \ln(w_{j}))}}{\sum_{k=1}^{J} \exp{(\epsilon \ln(w_{k}))}}L
\end{equation}
with $L$ the total labor force.

\item To simplify the analyses and abstract from strategic interactions in wage setting, assume that the number of firms $J$ is very large such that we can write:
\begin{equation} \label{eq_firm_level_labor_supply}
    \ln(l_{j}) = \epsilon \ln(w_{j}) + \ln(\lambda L)
\end{equation}
where $\lambda$ is a constant common across all firms. 

\item Interpret equation (\ref{eq_firm_level_labor_supply}). What does it mean if $ \epsilon \rightarrow 0 $?

\end{itemize}

\subsubsection{Firm optimization}

\begin{itemize}

\item Assume that firms have the following production functions:
\begin{equation}
    q_{j} = \psi_{j} \ln(l_{j})
\end{equation}
where $\psi_{j}$ is a firm-specific productivity shifter and the marginal product of labor is decreasing in employment. Both turn out to play an important role in our comparative statics later. Firms maximize profits by posting a wage that minimizes labor costs given equation (\ref{eq_firm_level_labor_supply}):
\begin{equation}
    \max_{w_{j}} \left[ \psi_{j} \ln(l_{j}) - w_{j} l_{j} \right] \text{\hspace{0.2cm} s.t. \hspace{0.2cm}} l_{j}=l_{j}(w_{j})
\end{equation}
Note that we assume employers do not observe the idiosyncratic taste for amenities of any given worker given by equation (\ref{eq_worker_preferences}). This information asymmetry implies employers cannot price discriminate between workers. Instead, if a firm wants to hire more workers it needs to offer higher wages to all workers.

\item Using the first-order conditions, derive the following expression:
\begin{equation} \label{eq_foc_log_wage}
   \ln(w_{j}) = \ln \left( \psi_{j} \frac{\epsilon}{1+ \epsilon} \right) - \ln(l_{j})
\end{equation}

\end{itemize}

\subsubsection{Equilibrium}

\begin{itemize}

\item Given equations (\ref{eq_firm_level_labor_supply}) and (\ref{eq_foc_log_wage}), derive expressions for firm-level wages and employment in equilibrium by replacing the question marks in the following expressions:
\begin{equation} \label{eq_log_wage}
   \ln(w_{j}) = \textbf{?} - \frac{1 }{1+ \epsilon} \ln(\lambda L)
\end{equation}
where the right-hand side only depends on the model's parameters. Equilibrium firm-level employment is given by:
\begin{equation} \label{eq_log_employment}
   \ln(l_{j}) = \textbf{?} + \frac{1}{1+ \epsilon} \ln(\lambda L)
\end{equation}

\item Could an increased difference in firm-level productivity over time explain the rising importance of firm effects and the increased sorting of workers into high-wage firms?

\end{itemize}

\end{document}